Using the copy tensor, we can write this as:
\begin{align}\label{eqn:parametrized-prob}
    \PP(X_1,\cdots,X_4) \simeq
    \begin{tikzpicture}[baseline=\baseline]
    \tikzset{tensor/.style = {minimum size = 0.5cm,shape = circle,thick,draw=black,inner sep = 0pt}, edge/.style   = {thick,line width=.4mm},every loop/.style={}}
    \def\y{-0.35}
    \def\x{0}
  \node[tensor,fill=red!30!white] (D) at (\x-1,\y-0.25){$\Gt_1$};
  \node[tensor,fill=red!50!white] (A) at (\x,\y-0.25){\scalebox{0.85}{$\Gt_2$}};
  \node[tensor,fill=blue!50!white] (B) at (\x+1,\y-0.25){$\Gt_3$};
   \node[tensor,fill=blue!20!white] (C) at (\x+2,\y-0.25){$\Gt_4$};
   \node[tensor,fill=red!30!white] (G1) at (\x-1,\y+1){$\Gt_1$};
  \node[tensor,fill=red!50!white] (G2) at (\x,\y+1){\scalebox{0.85}{$\Gt_2$}};
  \node[tensor,fill=blue!50!white] (G3) at (\x+1,\y+1){$\Gt_3$};
   \node[tensor,fill=blue!20!white] (G4) at (\x+2,\y+1){$\Gt_4$};
    \draw[edge](G1) -- (G2)node [midway,above] {\edgelabel{$R_1$}};
    \draw[edge](G2) -- (G3)node [midway,above] {\edgelabel{$R_2$}};
    \draw[edge](G3) -- (G4)node [midway,above] {\edgelabel{$R_3$}};
    \draw[edge](D) -- (A) node [midway,above] {\edgelabel{$R_1$}};
    \draw[edge](A) -- (B) node [midway,above] {\edgelabel{$R_2$}};
    \draw[edge](B) -- (C) node [midway,above] {\edgelabel{$R_3$}};
    \draw  node[fill = black,circle,inner sep=0pt,minimum size=4pt] (m_1) at (\x-1,\y+0.35) {};
    \draw  node[fill = black,circle,inner sep=0pt,minimum size=4pt] (m_2) at (\x,\y+0.35) {};
    \draw  node[fill = black,circle,inner sep=0pt,minimum size=4pt] (m_3) at (\x+1,\y+0.35) {};
    \draw  node[fill = black,circle,inner sep=0pt,minimum size=4pt] (m_4) at (\x+2,\y+0.35) {};
    \draw[edge](D)--(G1);
    \draw[edge](A)--(G2);
    \draw[edge](B)--(G3);    \draw[edge](C)--(G4);
    \draw[edge](m_1) -- (\x-1.35,\y+0.35);
    \draw[edge](m_4) -- (\x+2.35,\y+0.35);
    \draw[edge](m_2) -- (\x-0.35,\y+0.35);
    \draw[edge](m_3) -- (\x+1.35,\y+0.35);
\end{tikzpicture}
\end{align}

To ensure that the probabilities sum to one, we can choose to model un-normalized squared-root probabilities, since, as we will see, computing the normalization constant is very easy when the probability tensor is in the TT format. We thus assume that
$$\PP(X_1=i_1,\cdots,X_N=i_N) = \frac{(\Tt_{i_1,\cdots,i_N})^2}{\zeta},$$
where $\zeta = \sum_{i_1,\cdots,i_N}(\Tt_{i_1,\cdots,i_N})^2$ is the normalization constant. Similarly to efficient operations in the TT format presented in Section~\ref{sec:efficient-operations}, this normalization factor can be efficiently computed in the TT format: 
$$\zeta= \sum_{i_1,\cdots,i_4}\Tt_{i_1,\cdots,i_4}^2
=
\begin{tikzpicture}[baseline=\baseline]
    \tikzset{tensor/.style = {minimum size = 0.5cm,shape = circle,thick,draw=black,inner sep = 0pt}, edge/.style   = {thick,line width=.4mm},every loop/.style={}}
    \def\y{-0.35}
    \def\x{0}
  \node[tensor,fill=red!30!white] (D) at (\x-1,\y-0.25){$\Gt_1$};
  \node[tensor,fill=red!50!white] (A) at (\x,\y-0.25){\scalebox{0.85}{$\Gt_2$}};
    \node[tensor,fill=blue!50!white] (B) at (\x+1,\y-0.25){$\Gt_3$};
    \node[tensor,fill=blue!20!white] (C) at (\x+2,\y-0.25){$\Gt_4$};
     \node[tensor,fill=red!30!white] (G1) at (\x-1,\y+1){$\Gt_1$};
    \node[tensor,fill=red!50!white] (G2) at (\x,\y+1){\scalebox{0.85}{$\Gt_2$}};
    \node[tensor,fill=blue!50!white] (G3) at (\x+1,\y+1){$\Gt_3$};
    \node[tensor,fill=blue!20!white] (G4) at (\x+2,\y+1){$\Gt_4$};

    \draw[edge](G1) -- (G2)node [midway,above] {\edgelabel{$R_1$}};
    \draw[edge](G2) -- (G3)node [midway,above] {\edgelabel{$R_2$}};
    \draw[edge](G3) -- (G4)node [midway,above] {\edgelabel{$R_3$}};
    \draw[edge](D) -- (A) node [midway,above] {\edgelabel{$R_1$}};
    \draw[edge](A) -- (B) node [midway,above] {\edgelabel{$R_2$}};
    \draw[edge](B) -- (C) node [midway,above] {\edgelabel{$R_3$}};
    \draw[edge](G1)--(D);
    \draw[edge](G2)--(A);
    \draw[edge](G3)--(B);
    \draw[edge](G4)--(C);
\end{tikzpicture}.
$$
The idea above is also called \emph{Born Machine}~\citep{glasser2019expressive,han2018unsupervised}.
Similarly to what we showed in Section~\ref{sub:probs}, we can compute marginal distributions of TT parametrized probability distributions~(see Eq.~\eqref{eqn:parametrized-prob}) as tensor networks:
\begin{enumerate}
    \item %
    Using the law of total probabilities, the marginal distribution of $X_2,X_3$, and $X_4$ can be expressed as 
    \begin{align*}
    \PP(X_2=i_2, X_3=i_3, X_4=i_4) 
    &= \sum_{i_1=1}^{d_1}\PP(X_1=i_1,\cdots,X_4=i_4)\\ 
    &=
    \begin{tikzpicture}[baseline=\baseline]
    \tikzset{tensor/.style = {minimum size = 0.5cm,shape = circle,thick,draw=black,inner sep = 0pt}, edge/.style   = {thick,line width=.4mm},every loop/.style={}}
    \def\y{-0.35}
    \def\x{0}
  \node[tensor,fill=red!30!white] (D) at (\x-1,\y-0.25){$\Gt_1$};
  \node[tensor,fill=red!50!white] (A) at (\x,\y-0.25){\scalebox{0.85}{$\Gt_2$}};
  \node[tensor,fill=blue!50!white] (B) at (\x+1,\y-0.25){$\Gt_3$};
   \node[tensor,fill=blue!20!white] (C) at (\x+2,\y-0.25){$\Gt_4$};
   \node[tensor,fill=red!30!white] (G1) at (\x-1,\y+1){$\Gt_1$};
  \node[tensor,fill=red!50!white] (G2) at (\x,\y+1){\scalebox{0.85}{$\Gt_2$}};
  \node[tensor,fill=blue!50!white] (G3) at (\x+1,\y+1){$\Gt_3$};
   \node[tensor,fill=blue!20!white] (G4) at (\x+2,\y+1){$\Gt_4$};
    \draw[edge](G1) -- (G2)node [midway,above] {\edgelabel{$R_1$}};
    \draw[edge](G2) -- (G3)node [midway,above] {\edgelabel{$R_2$}};
    \draw[edge](G3) -- (G4)node [midway,above] {\edgelabel{$R_3$}};
    \draw[edge](D) -- (A) node [midway,above] {\edgelabel{$R_1$}};
    \draw[edge](A) -- (B) node [midway,above] {\edgelabel{$R_2$}};
    \draw[edge](B) -- (C) node [midway,above] {\edgelabel{$R_3$}};
    \draw  node[fill = black,circle,inner sep=0pt,minimum size=4pt] (m_1) at (\x-1,\y+0.35) {};
    \draw  node[fill = black,circle,inner sep=0pt,minimum size=4pt] (m_2) at (\x,\y+0.35) {};
    \draw  node[fill = black,circle,inner sep=0pt,minimum size=4pt] (m_3) at (\x+1,\y+0.35) {};
    \draw  node[fill = black,circle,inner sep=0pt,minimum size=4pt] (m_4) at (\x+2,\y+0.35) {};
    \draw[edge](D)--(G1);
    \draw[edge](A)--(G2);
    \draw[edge](B)--(G3);    \draw[edge](C)--(G4);
    \draw[edge](m_1) -- (\x-1.35,\y+0.35);
    \draw  node[fill = black,circle,inner sep=0pt,minimum size=4pt] (m) at (\x-1.35,\y+0.35) {};
    \draw[edge](m_4) -- (\x+2.35,\y+0.35) node [right] {\scalebox{0.7}{\indcolor{$i_4$}}};;
    \draw[edge](m_2) -- (\x-0.35,\y+0.35) node [left] {\scalebox{0.7}{\indcolor{$i_2$}}};;
    \draw[edge](m_3) -- (\x+1.35,\y+0.35) node [right] {\scalebox{0.7}{\indcolor{$i_3$}}};;
\end{tikzpicture}
=
\sum_{i_1}^{d_1}\Tt_{i_1,\cdots,i_4},
    \end{align*}
for all $i_2\in[d_1],i_3\in[d_3]$ and $i_4\in[d_4]$. 
\item We can similarly express the remaining marginal distributions. For example, the marginal distribution of $X_2$ can be written as
\begin{align*}
    \PP(X_2=i_2) 
    &= \sum_{i_1,i_3,i_4=1}^{d_1,d_3,d_4}\PP(X_1=i_1,\cdots,X_4=i_4)\\ 
    &=
    \begin{tikzpicture}[baseline=\baseline]
    \tikzset{tensor/.style = {minimum size = 0.5cm,shape = circle,thick,draw=black,inner sep = 0pt}, edge/.style   = {thick,line width=.4mm},every loop/.style={}}
    \def\y{-0.35}
    \def\x{0}
  \node[tensor,fill=red!30!white] (D) at (\x-1,\y-0.25){$\Gt_1$};
  \node[tensor,fill=red!50!white] (A) at (\x,\y-0.25){\scalebox{0.85}{$\Gt_2$}};
  \node[tensor,fill=blue!50!white] (B) at (\x+1,\y-0.25){$\Gt_3$};
   \node[tensor,fill=blue!20!white] (C) at (\x+2,\y-0.25){$\Gt_4$};
   \node[tensor,fill=red!30!white] (G1) at (\x-1,\y+1){$\Gt_1$};
  \node[tensor,fill=red!50!white] (G2) at (\x,\y+1){\scalebox{0.85}{$\Gt_2$}};
  \node[tensor,fill=blue!50!white] (G3) at (\x+1,\y+1){$\Gt_3$};
   \node[tensor,fill=blue!20!white] (G4) at (\x+2,\y+1){$\Gt_4$};
    \draw[edge](G1) -- (G2)node [midway,above] {\edgelabel{$R_1$}};
    \draw[edge](G2) -- (G3)node [midway,above] {\edgelabel{$R_2$}};
    \draw[edge](G3) -- (G4)node [midway,above] {\edgelabel{$R_3$}};
    \draw[edge](D) -- (A) node [midway,above] {\edgelabel{$R_1$}};
    \draw[edge](A) -- (B) node [midway,above] {\edgelabel{$R_2$}};
    \draw[edge](B) -- (C) node [midway,above] {\edgelabel{$R_3$}};
    \draw  node[fill = black,circle,inner sep=0pt,minimum size=4pt] (m_1) at (\x-1,\y+0.35) {};
    \draw  node[fill = black,circle,inner sep=0pt,minimum size=4pt] (m_2) at (\x,\y+0.35) {};
    \draw  node[fill = black,circle,inner sep=0pt,minimum size=4pt] (m_3) at (\x+1,\y+0.35) {};
    \draw  node[fill = black,circle,inner sep=0pt,minimum size=4pt] (m_4) at (\x+2,\y+0.35) {};
    \draw[edge](D)--(G1);
    \draw[edge](A)--(G2);
    \draw[edge](B)--(G3);    
    \draw[edge](C)--(G4);
    \draw[edge](m_1) -- (\x-1.35,\y+0.35);
    \draw  node[fill = black,circle,inner sep=0pt,minimum size=4pt] (m_5) at (\x-1.35,\y+0.35) {};
    \draw[edge](m_4) -- (\x+2.35,\y+0.35);
    \draw  node[fill = black,circle,inner sep=0pt,minimum size=4pt] (m_6) at (\x+2.35,\y+0.35) {};
    \draw[edge](m_2) -- (\x-0.35,\y+0.35) node [left] {\scalebox{0.7}{\indcolor{$i_2$}}};;
    \draw[edge](m_3) -- (\x+1.35,\y+0.35);
    \draw  node[fill = black,circle,inner sep=0pt,minimum size=4pt] (m_7) at (\x+1.35,\y+0.35) {};
\end{tikzpicture}
=
\sum_{i_1,i_3,i_4=1}^{d_1,d_3,d_4}\Tt_{i_1,\cdots,i_4},
    \end{align*}
\end{enumerate}
\subsection{Computing Expectations of Random Tensor Networks} 
We will now show how the tensor networks formalism can help us derive complex properties of random tensors. 
We begin this section by presenting two very useful propositions.
The first one is a trivial consequence of linearity of expectation, but will prove to be very useful in this section.
\begin{proposition}\label{prop:inner-expectation}
    Let $\At$ and $\Bt$ be two independent random tensor networks. tensor networks. Then their inner product in expectation is 
$$\EE\left(\begin{tikzpicture}[baseline=\baseline]
    \tikzset{tensor/.style = {minimum size = 0.5cm,shape = circle,thick,draw=black,inner sep = 0pt}, edge/.style = {thick,line width=.4mm},every loop/.style={}}
    \def\y{0}
    \def\x{0}
    \node[tensor,fill=red!50!blue!50!white] (At) at (\x,\y){\scalebox{0.85}{$\At$}};
    \node[tensor,fill=red!30!blue!50!white] (Bt) at (\x+1,\y){\scalebox{0.85}{$\Bt$}};
    \draw[edge] (At) -- (Bt);
    \draw[edge] (At) to [out=45,in=135] (Bt);
    \draw[edge] (At) to [out=-45,in=-135] (Bt);
    \draw[edge](At)--(\x-0.85,\y);
    \draw[edge](At)--(\x-0.75,\y+0.5);    \draw[edge](At)--(\x-0.75,\y-0.5);
    \draw[edge](Bt)--(\x+1.85,\y);
    \draw[edge](Bt)--(\x+1.75,\y+0.5);    \draw[edge](Bt)--(\x+1.75,\y-0.5);
\end{tikzpicture} \right)
=
\begin{tikzpicture}[baseline=\baseline]
    \tikzset{tensor/.style = {minimum size = 0.8cm,shape = circle,thick,draw=black,inner sep = 0pt}, edge/.style = {thick,line width=.4mm},every loop/.style={}}
    \def\y{0}
    \def\x{0}
    \node[tensor,fill=red!50!blue!50!white] (At) at (\x,\y){\scalebox{0.85}{$\EE(\At)$}};
    \node[tensor,fill=red!30!blue!50!white] (Bt) at (\x+1.5,\y){\scalebox{0.85}{$\EE(\Bt)$}};
    \draw[edge] (At) -- (Bt);
    \draw[edge] (At) to [out=45,in=135] (Bt);
    \draw[edge] (At) to [out=-45,in=-135] (Bt);
    \draw[edge](At)--(\x-0.85,\y);
    \draw[edge](At)--(\x-0.75,\y+0.5);    \draw[edge](At)--(\x-0.75,\y-0.5);
    \draw[edge](Bt)--(\x+2.35,\y);
    \draw[edge](Bt)--(\x+2.25,\y+0.5);    \draw[edge](Bt)--(\x+2.25,\y-0.5);
\end{tikzpicture}.
$$
\begin{proof}
The result directly follows from the linearity of the expected value and the independence of $\At$ and $\Bt$.
\end{proof}
\end{proposition}
The second one shows an interesting properties of matrices whose entries are drawn from a Gaussian distribution: the expectation of the outer product of such a matrix with itself reduces to a very simple tensor network. 
\begin{proposition}\label{prop:outer-prod}Let $\Ab\in\RR^{m\times n}$ be a random matrix whose elements are drawn identically and independently from the standard normal distribution. The expected value of the outer product of $\Ab$ with itself is 
$$\EE\left(
\begin{tikzpicture}[baseline=\baseline]
    \tikzset{tensor/.style = {minimum size = 0.5cm,shape = circle,thick,draw=black,inner sep = 0pt}, edge/.style = {thick,line width=.4mm},every loop/.style={}}
    \def\y{0.15}
    \def\x{0}
    \node[tensor,fill=green!50!blue!20!white] (A) at (\x,\y){\scalebox{0.85}{$\Ab$}};
    \node[tensor,fill=green!50!blue!20!white] (B) at (\x+1.5,\y){\scalebox{0.85}{$\Ab$}};
    \draw[edge](A) -- (\x-0.55,\y-0.55);
    \node[draw=none] () at (\x-0.5,\y-0.3) {\scalebox{0.7}{\dimcolor{$m$}}};
    \draw[edge](A) -- (\x+0.55,\y-0.55);
    \node[draw=none] () at (\x+0.5,\y-0.3) {\scalebox{0.7}{\dimcolor{$n$}}};
    \draw[edge](B) -- (\x+1,\y-0.55);
    \node[draw=none] () at (\x+1,\y-0.3) {\scalebox{0.7}{\dimcolor{$m$}}};
    \draw[edge](B) -- (\x+2,\y-0.55);
    \node[draw=none] () at (\x+2,\y-0.3) {\scalebox{0.7}{\dimcolor{$n$}}};
\end{tikzpicture}
\right)
=
\begin{tikzpicture}[baseline=\baseline]
    \tikzset{tensor/.style = {minimum size = 0.5cm,shape = circle,thick,draw=black,inner sep = 0pt}, edge/.style = {thick,line width=.4mm},every loop/.style={}}
    \def\y{0}
    \def\x{0}
    \draw[edge] (\x+0.15,\y-0.1) to [out=100,in=80] (\x+1.85,\y-0.1);
    \draw[edge] (\x+1.45,\y-0.1) to [out=100,in=80] (\x+3,\y-0.1);
    \node[draw=none] () at (\x+2.25,\y+0.5) {\scalebox{0.6}{\dimcolor{$n$}}};
    \node[draw=none] () at (\x+1,\y+0.5) {\scalebox{0.6}{\dimcolor{$m$}}};
\end{tikzpicture},
$$
where the right-hand side
is the outer product of two Kronecker delta~(i.e., $
\begin{tikzpicture}[baseline=\baseline]
    \tikzset{tensor/.style = {minimum size = 0.5cm,shape = circle,thick,draw=black,inner sep = 0pt}, edge/.style = {thick,line width=.4mm},every loop/.style={}}
    \def\y{0}
    \def\x{0}
    \draw[edge] (\x+0.15,\y-0.1) to [out=100,in=80] (\x+1.85,\y-0.1);
    \draw[edge] (\x+1.45,\y-0.1) to [out=100,in=80] (\x+3,\y-0.1);
    \node[draw=none] () at (\x+0.15,\y-0.25) {\scalebox{0.7}{\indcolor{$i_1$}}};
    \node[draw=none] () at (\x+3,\y-0.25) {\scalebox{0.7}{\indcolor{$i_4$}}};
    \node[draw=none] () at (\x+1.85,\y-0.25) {\scalebox{0.7}{\indcolor{$i_3$}}};
    \node[draw=none] () at (\x+1.45,\y-0.25) {\scalebox{0.7}{\indcolor{$i_2$}}};
\end{tikzpicture}
=\delta_{i_1i_3}\delta_{i_2i_4}$ for $i_1,i_3\in[m]$ and $i_2,i_4\in[n]$).
\begin{proof}
Since the entries of $\Ab$ are drawn independently from $\mathcal{N}(0,1)$, its vectorization follows a multivariate normal distribution, $\vectorize(\Ab) \sim \mathcal{N}(\mathbf{0},\Ib_{mn})$, and thus $\EE(\vectorize(\Ab)\vectorize(\Ab)^\ts) = \Ib_{mn}$. This means that 
$$ \EE(\Ab_{i_1i_2}\Ab_{i_3i_4})=
\begin{cases}
        1 & \text{ if } i_1=i_3\text{ and } i_2=i_4\\
        0 & \text{ otherwise, }
    \end{cases} 
    $$
hence, $\EE(\Ab\outprod\Ab)_{i_1i_2i_3i_4}=\EE(\Ab_{i_1i_2}\Ab_{i_3i_4}) = \delta_{i_1i_3}\delta_{i_2i_4}=
\begin{tikzpicture}[baseline=\baseline]
    \tikzset{tensor/.style = {minimum size = 0.5cm,shape = circle,thick,draw=black,inner sep = 0pt}, edge/.style = {thick,line width=.4mm},every loop/.style={}}
    \def\y{0}
    \def\x{0}
    \draw[edge] (\x+0.15,\y-0.1) to [out=100,in=80] (\x+1.85,\y-0.1);
    \draw[edge] (\x+1.45,\y-0.1) to [out=100,in=80] (\x+3,\y-0.1);
    \node[draw=none] () at (\x+0.15,\y-0.25) {\scalebox{0.7}{\indcolor{$i_1$}}};
    \node[draw=none] () at (\x+3,\y-0.25) {\scalebox{0.7}{\indcolor{$i_4$}}};
    \node[draw=none] () at (\x+1.85,\y-0.25) {\scalebox{0.7}{\indcolor{$i_3$}}};
    \node[draw=none] () at (\x+1.45,\y-0.25) {\scalebox{0.7}{\indcolor{$i_2$}}};
\end{tikzpicture}$, for all $i_1,i_3\in[m]$ and $i_2,i_4\in[n]$.

